\documentclass{article}
\usepackage[utf8]{inputenc}

% Tutaj można zmienić grubość marginesu.
\usepackage[left=2cm, top=2cm, right=2cm, bottom=1cm, includeheadfoot, headheight=50pt, a4paper]{geometry}
\usepackage{polski}
\usepackage{fancyhdr}
\usepackage{lastpage}
\usepackage{amsmath}
\usepackage{amssymb}
\usepackage{wrapfig}
\usepackage{soul}

%%%%%%%%%%%%%%%%%%%%%%%%%%%%
%%% Możesz dodać pakiety %%%
%%%%%%%%%%%%%%%%%%%%%%%%%%%%

\usepackage{amsfonts}
\usepackage{amsmath}
\usepackage{enumitem}
\usepackage{braket}
\usepackage{amssymb}
\usepackage{amsthm}
\usepackage{xcolor}

\usepackage{hyperref}
\usepackage{graphicx}
\usepackage{mathtools}
\usepackage{svg}
\usepackage{listings}

\newtheorem{theorem}{Theorem}[section]
\newtheorem{lemma}{Lemat}

\newcommand{\imieNazwisko}{Dominika Piętka}
\newcommand{\nrpd}{1}

% Tutaj jest numeracja stron
\pagestyle{fancy}

\fancyhf{}

\cfoot{
Strona $\thepage$ z $\pageref{LastPage}$
}

\renewcommand{\headrulewidth}{0pt}


\let\leq=\leqslant % redefine
\let\geq=\geqslant % redefine


\renewcommand\qedsymbol{$\blacksquare$}

\title{Praca domowa \nrpd}
\author{Dominika Piętka}
\date{6 października 2024}



\begin{document}
\fontsize{12}{16}\selectfont



\maketitle

\textbf{Zadanie 1}
$$\{ x \in \mathbb{Z} \mid \exists y \in \mathbb{Z} (y > x \land \forall x \in \mathbb{Z} (x > y \to x > 1.1y)) \}$$ 
  Niech $Y = \{ y \in \mathbb{Z} \mid  \forall x \in \mathbb{Z} (x > y \to x > 1.1y)) \}$. Gdy $x \leq y$ implikacja zawsze jest prawdziwa, rozważmy więc $x > y$. Skoro $x, y \in \mathbb{Z}$ to $x - y \geq 1$. Ponadto, $x > 1.1 y \iff x - y > 0.1y$. $x - y \geq 1$, zatem musi być spełnione $0.1y < 1 \iff y < 10$. W takim razie $Y = (-\infty, 10)$. 

  $\{ x \in \mathbb{Z} \mid \exists y \in \mathbb{Z} (y > x \land \forall x \in \mathbb{Z} (x > y \to x > 1.1y)) \} \iff \{ x \in \mathbb{Z} \mid \exists y \in Y (y > x) \} \iff \newline \iff \{ x\in \mathbb{Z} \mid x < 9 \}$, co dla liczb całkowitych jest równoważne z $\{ x\in \mathbb{Z} \mid x \leq 8 \} $:

  $$ \includegraphics[width=0.7\textwidth]{Pmat 1.1.png}$$

  \textbf{Zadanie 2}
  \begin{enumerate}
      % \item Założeniem dowod nie wprost jest zaprzeczenie tezy, a $\neq (p \land q) = \neq p \lor \neq q$, zatem założenie powinno brzmieć "ciąg \st{rosnący} niemalejący \st{i} lub nieograniczony"
      \item Dowód nie wprost polega na zaprzeczeniu tezie i wykazaniu, że prowadzi to do sprzeczności. Teza ma postać $(p \land q) \rightarrow r$, zatem $\neg ((p \land q) \rightarrow r) \leftrightarrow (p \land q) \land r$. Dowód powinien więc brzmieć \textit{Przeprowadzimy dowód nie wprost. Załóżmy, że pewien ciąg
jest malejący, ograniczony i rozbieżny.}
    \item Przypuszczalnie \textit{ciąg rosnący i nieograniczony} miał być zaprzeczeniem \textit{ciągu malejącego i ograniczonego}, co jest błędne z dwóch powodów:
    \begin{enumerate}
        \item zaprzeczeniem \textit{malejący} jest \textit{niemalejący}
        \item $\neg (p \land q) \leftrightarrow \neg q \lor \neg q$, zatem zdanie \textit{Załóżmy, że pewien ciąg jest rosnący i nieograniczony} powinno brzmieć \textit{Załóżmy, że pewien ciąg jest niemalejący lub nieograniczony}.
    \end{enumerate}
  \end{enumerate}


\end{document}
